\section{Train-on-test inflation of a published VLA-RL result}
\label{sec:inflation}

LIBERO is tested on 50 fixed initial states per task. The released SimpleVLA-RL LIBERO-10 checkpoint was trained by RL on exactly those states. We measured how much of its published number this accounts for, as a preregistered secondary measurement: the RL checkpoint and its own SFT starting point were each evaluated greedily on the 500 official episodes and on 500 placements freshly drawn by LIBERO's own sampler (the same distribution, not a perturbation). Inflation is the difference of the two official--fresh gaps, which removes any difficulty difference between the placement sets; the CI is a bootstrap over episodes.

\begin{table}[t]
  \centering
  \small
  \caption{\textbf{Train-on-test inflation} (LIBERO-10, success \%).}
  \label{tab:inflation}
  \begin{tabular}{@{}lccc@{}}
    \toprule
    checkpoint & official 500 & fresh 500 & gap \\
    \midrule
    SFT start (never saw the test states) & 14.6 & 20.0 & $-5.4$ \\
    RL, trained on the official states & 87.4 & 76.0 & $+11.4$ \\
    \midrule
    inflation & & & $+16.8$ $[9.8, 23.4]$ \\
    \bottomrule
  \end{tabular}
\end{table}

The inflation is $+16.8$ points $[+9.8, +23.4]$ (\cref{tab:inflation}); a replication on 500 further never-seen placements gives $+16.0$ $[+9.2, +22.4]$, and the loss is concentrated on two of the ten tasks. The checkpoint is still far better than its start on fresh placements ($+56$ points), so the method works; the point is that about a fifth of the headline number is memorisation of the test states. Every number in this paper is measured on the official states \emph{without} ever training, sampling or selecting on them, which is why our finals should be compared with the ``fresh'' column of such results rather than the ``official'' one.
